\section{A string check both finds and labels delayed information needs}
\label{sec:needs}
A step $t$ has a \emph{delayed information need} when the text the agent types at $t$ appears verbatim in the OCR text of an earlier screen $i<t$ (a \emph{source screen}) and does not appear on screen $t$. The gap is $t-i$ for the latest such $i$. The second clause matters: when the text is on the current screen the agent needs no memory, and keeping such steps rewards any method that reads recent screens (\S\ref{sec:limits}).

\paragraph{The verifier mines needs and labels screens in one pass.} A candidate is any TYPE action whose argument has three or more characters and does not occur in the task description; we normalize strings by lower-casing and stripping whitespace and punctuation. We run EasyOCR on every screenshot of the sampled trajectories and test whether the argument is a substring of an earlier screen's OCR text and absent from the current screen's. In the first sample of 300 GUI-Odyssey trajectories, 127 of 431 candidates pass the OCR check (269 fail, most of them OCR misreads on inspection, and 35 lack OCR output). Removing 36 with the text on the current screen leaves 91 verified cases; removing 6 whose text appears in the task's step-by-step instruction field leaves 85. The same test, applied to every screen instead of stopping at the first hit, returns the set of screens that contain the text. That set is the label: screen $i$ is positive if it is in the set. The count of needs is a lower bound, and the rule covers only verbatim re-typing.

\begin{table}[t]
\centering\scriptsize
\caption{AITZ's 7.4-step tasks contain no delayed information need; GUI-Odyssey's 15.4-step tasks contain one in 31\% of sampled trajectories. We mine only the sampled trajectories.}
\label{tab:prev}
\begin{tabular}{lrrrl}
\toprule
 & trajectories & steps & mean length & verified delayed needs \\
\midrule
AITZ~\citep{aitz} & 2,504 & 18,643 & 7.4 & 0 \\
GUI-Odyssey~\citep{guiodyssey} (full) & 7,735 & 118,791 & 15.4 & not mined \\
\quad sampled, complete OCR & 273 & 4,338 & 15.9 & 89 steps (2.05\%); 31\% of trajectories hold at least one \\
\bottomrule
\end{tabular}
\end{table}

\paragraph{The need is common, far, and inconspicuous.} Table~\ref{tab:prev} gives prevalence from the 91 verified cases before the instruction-field filter, 89 of which fall in the 273 sampled trajectories with complete OCR coverage. The gap has median 5 and maximum 27 screens, 56\% of gaps are 5 or more, and on its source screen the needed text sits at a median rank of 13 among 33 OCR elements and is never first.

\paragraph{Four samples give 188 training cases and 93 external ones.} The \emph{short set} holds the 85 cases above (the need's step index $t$ has median 11). The \emph{long set} holds 53 cases with $t\ge16$ from a second sample of 300 trajectories of 20 or more steps (127 before the leakage filter, all 74 removals being leakage). The \emph{late-mined set} holds 50 cases with $t\ge16$ from a third sample of 200 trajectories. These three give 188 cases in 170 trajectories, and all training and cross-validation use them. The \emph{external set} holds 93 cases from a fourth sample of 400 trajectories of 16 or more steps. We mined it after fixing every design decision and evaluate on it once (\S\ref{sec:ext}); its needs sit mostly at steps 12--15 (72 of 93).
